\documentclass[10pt,a4paper]{amsart}
\usepackage[margin=19mm]{geometry}
\usepackage{amsmath,amssymb,mathtools}
\usepackage[scr=rsfs,cal=euler]{mathalfa}
\usepackage{xfrac}
\usepackage[unicode,hidelinks,bookmarks=false]{hyperref}
\newcommand{\ie}{i.e.\ }
\newcommand{\cf}{cf.\ }
\newcommand{\resp}{respectively\ }
\newcommand{\Subsection}{\subsection}
\providecommand{\nobreakdash}{\nobreak\hbox{-}}
\setlength{\emergencystretch}{2em}
\multlinegap0pt
\usepackage{bm}
\usepackage{bbm}
\usepackage{dsfont}
\usepackage{xspace}
\usepackage{cancel}
\usepackage{mathtools}
\usepackage{cleveref}
\usepackage{amsthm}
\makeatletter
\@ifundefined{theorem}{\newtheorem{theorem}{Theorem}}{}
\@ifundefined{proposition}{\newtheorem{proposition}[theorem]{Proposition}}{}
\makeatother
\newcommand*{\R}{\mathbb{R}}                                    % Real numbers
\newcommand*{\cB}{\mathcal{B}}
\newcommand*{\eqdef}{\stackrel{\mbox{\normalfont\tiny def}}{=}} % definition by equality                                      
\newcommand*{\vF}{\vec{F}}
\newcommand*{\veps}{\varepsilon}                                % Nice-looking epsilon
\newcommand*{\vn}{\vec{n}}
\title[A dichotomy for hypergraph Zarankiewicz problems on axis-par]{A dichotomy for hypergraph Zarankiewicz problems on axis-parallel boxes}
\author{Ting-Wei Chao, Zichao Dong, Hong Liu, Xichao Shu, Shuaichao Wang}
\date{Source: arXiv:2604.20815v1 (CC BY 4.0)}
\begin{document}
\maketitle

\noindent\emph{Source and licence.} A dichotomy for hypergraph Zarankiewicz problems on axis-parallel boxes. Version arXiv:2604.20815v1 (CC BY 4.0), distributed under the Creative Commons Attribution 4.0 International licence, as recorded in the arXiv licence metadata for this exact version and shown on the abstract page https://arxiv.org/abs/2604.20815v1. Full rights notice: licenses/arxiv-2604.20815-CC-BY-4.0.txt.\par\smallskip

\noindent\textit{Editorial scope.} Counted unit: the Proposition whose source label is \texttt{prop:upper\_canonical} in the article (source0.tex lines 905--908, source label \texttt{prop:upper\_canonical}), delivered together with the article's own complete original proof of it (source0.tex lines 912--935, one \texttt{proof} environment of 24 lines). No mathematics is written, repaired or re-derived: statement and proof are the article's own text line for line. The proof is detached from the statement in the source: the article states the result at lines 905--908 and prints its proof at lines 912--935, and the proof environment's own header names the statement, so the association is the article's own. Required context retained uncounted: prop:upper\_separated (lines 750--752). Every definition, display and label of those lines is retained unchanged. Omitted from this package and not redistributed: 1--749, 753--904, 909--911, 936--1253. Nothing inside a retained excerpt is omitted or altered: every retained body is delivered exactly as the article's lines are written, so no omission receipt of any kind accompanies this package. Citations: the retained excerpts carry no citation command at all, so no bibliography entry of the article is transcribed and no reference is redistributed. Reference adaptation: none was needed and none was made; every reference inside the delivered unit resolves inside it, so no unresolved reference or citation is produced. Printed slips kept literal: no printed word, symbol, hypothesis or number of the counted statement or of its proof was altered. Layout and class: the article's own class is \texttt{article} with packages including hyperref, bookmark, amsfonts, amsmath, amssymb, amsthm, enumitem, pgf, tikz, graphicx, bm, cancel, comment, etoolbox; the wrapper is the lane's pinned \texttt{amsart} editorial wrapper at 10pt on A4, which restates the theorem-like environments the retained text needs and adds the article's own macro definitions that the retained text needs (\texttt{\textbackslash R}, \texttt{\textbackslash cB}, \texttt{\textbackslash eqdef}, \texttt{\textbackslash vF}, \texttt{\textbackslash veps}, \texttt{\textbackslash vn}), with extra wrapper packages (bm, bbm, dsfont, xspace, cancel, mathtools, cleveref). The delivered numbering is the unit's own. Assets: no retained span contains a figure environment, a graphics-inclusion command, a PDF-page inclusion, a table or a TikZ picture; no image file is copied into this package, the package declares no asset dependency and no image file is redistributed with it. Licence: version arXiv:2604.20815v1 of this article is distributed under the Creative Commons Attribution 4.0 International licence, as recorded in the arXiv licence metadata for this exact version and shown on the abstract page https://arxiv.org/abs/2604.20815v1. A full rights notice is delivered at licenses/arxiv-2604.20815-CC-BY-4.0.txt. Characters: every retained excerpt is pure ASCII, so the delivered engine, XeLaTeX with its default Latin Modern text font, reports no missing character.
\begin{proposition} \label{prop:upper_separated}
    Suppose $d \ge 2$ and $t \ge 2$. There exists $M_d > 0$, a constant depending only on $d$, such that $z_d^*(\vn; t) \le M_d \cdot g_t(\vn)$ holds for every $d$-size-vector $\vn$. 
\end{proposition}

\begin{proposition} \label{prop:upper_canonical}
    Suppose $d \ge 2$ and $t \ge 2$. Then $z_d(\vn; t) = z_d^*(\vn; t) \le M_d \cdot g_t(\vn)$ holds for every $d$-size-vector $\vn$, where the constant $M_d > 0$ is given by \Cref{prop:upper_separated}.
    \vspace{-0.5em}
\end{proposition}

\begin{proof}[Proof of \Cref{prop:upper_canonical}]
    Fix an arbitrary canonical $\vF_d$-family $(\cB_1, \dots, \cB_d)$. For $\bm{b}_1 \in \cB_1, \dots, \bm{b}_d \in \cB_d$, if they intersect, then the $d$-direction-vector $\vF_d = \bigl( [d] \setminus \{1\}, \dots, [d] \setminus \{d\} \bigr)$ shows that the intersection must be a single point. For each $i \in [d]$ and any $\bm{b} \in \cB_i$, we consider the following enlargement of $\bm{b}$.
    \[
    \begin{array}{ccccccccccccccc}
        \bm{b} &\eqdef &{\scriptstyle [a_1, b_1]} &{\scriptstyle \times} &{\scriptstyle \cdots} &{\scriptstyle \times} &{\scriptstyle [a_{i-1}, b_{i-1}]} &{\scriptstyle \times} &{\scriptstyle \{a_i\}} &{\scriptstyle \times} &{\scriptstyle [a_{i+1}, b_{i+1}]} &{\scriptstyle \times} &{\scriptstyle \cdots} &{\scriptstyle \times} &{\scriptstyle [a_d, b_d]} \\[-1.25em]
        \rotatebox{-90}{$\mapsto$} & & & & & & & & & & & & & & \\[0.5em]
        \bm{b}^{\veps, i} &\eqdef &{\scriptstyle [a_1-\veps, b_1+\veps]} &{\scriptstyle \times} &{\scriptstyle \cdots} &{\scriptstyle \times} &{\scriptstyle [a_{i-1}-\veps, b_{i-1}+\veps]} &{\scriptstyle \times} &{\scriptstyle \{a_i\}} &{\scriptstyle \times} &{\scriptstyle [a_{i+1}-\veps, b_{i+1}+\veps]} &{\scriptstyle \times} &{\scriptstyle \cdots} &{\scriptstyle \times} &{\scriptstyle [a_d-\veps, b_d+\veps]}. 
    \end{array}
    \]
    Intersections are preserved under the enlargement. Moreover, if $\bm{b}_1 \in \cB_1, \dots, \bm{b}_d \in \cB_d$ do not intersect, then they still do not intersect after enlargement, given $\veps > 0$ is small enough. Thus, the intersection hypergraph of $(\cB_1^{\veps, 1}, \dots, \cB_d^{\veps, d})$ coincides with that of $(\cB_1, \dots, \cB_d)$, where $\cB_i^{\veps, i} \eqdef \{\bm{b}^{\veps, i} \mid \bm{b} \in \cB_i\}$. 

    After the enlargement, every intersection point $\bm{b}^{\veps, 1} \cap \dots \cap \bm{b}^{\veps, d} \, (\bm{b}^{\veps, i} \in \cB_i^{\veps, i})$ lies in the interior of the box $\bm{b}^{\veps, i}$ for each $i \in [d]$. So, we may shift each box $\bm{b} \in \cB_i^{\veps, i}$ along the $i$-th direction by a segment of length $\delta_{\bm{b}} \in \mathbb{R}$, which is far smaller than $\veps$, to separate all the boxes. Specifically, 
    \vspace{-0.5em}
    \begin{itemize}
        \item if $\bm{b}_1 \cap \dots \cap \bm{b}_d = \{p\}$, then $p + \delta_{\bm{b}_1} \vec{e}_1 + \dots + \delta_{\bm{b}_d} \vec{e}_d$ is the intersecting point of these boxes after shifting, where $\vec{e}_i$ denotes the $i$-th unit coordinate vector of $\R^d$ for $i \in [d]$; 
        \vspace{-0.5em}
        \item if $\bm{b}_1 \cap \dots \cap \bm{b}_d = \varnothing$, then after the shifting we still have $(\bm{b}_1 + \delta_{\bm{b}_1} \vec{e}_1) \cap \dots \cap (\bm{b}_d + \delta_{\bm{b}_d} \vec{e}_d) = \varnothing$, because all the boxes that we are considering in $\R^d$ are closed. 
    \end{itemize}
    \vspace{-0.5em}
    Therefore, we have separated the family $(\cB_1^{\veps,1}, \dots, \cB_d^{\veps,d})$ while preserving the intersection hypergraph. 

    We thus obtain $z_d(\vn; t) = z_d^*(\vn; t) \le M_d \cdot g_t(\vn)$, where ``$\le$'' follows from \Cref{prop:upper_separated}. 
    % Secondly, since every intersection in $(\cB_1', \dots, \cB_d')$ is normal, for each box in $\cB_1' \cup \dots \cup \cB_d'$ we can translate it slightly along any direction without changing the intersection hypergraph while keeping all intersections normal. Thus, we can arbitrarily order all boxes in $\cB_1' \cup \dots \cup \cB_d'$ and apply a tiny translation one by one to each of them (say $\bm{b}_*$) so that $S_i(\bm{b}_*)$ becomes disjoint with $S_i(\bm{b})$ for every $i \in [d]$ and every $\bm{b} \in \cB_1' \cup \dots \cup \cB_d' \setminus \{\bm{b}_*\}$. After applying all such small translations, we obtain a new separated family $(\widetilde{\cB}_1, \dots, \widetilde{\cB}_d)$ which preserves the intersection hypergraph of $(\cB_1, \dots, \cB_d)$. 
\end{proof}

\end{document}
